这本书想解决一个具体的困境：你能写 Python，也听说过 monad，看过几篇讲范畴论的博客，
但这两件事在你脑子里始终是分开的------数学是数学，代码是代码。

我的写法是\textbf{从能跑的代码倒推数学}，而不是从定义推向代码。每一章都先有一个跑不起来、
或者写起来极其难看的程序，再由它逼出一个概念。第 7 章的函子是从''给一堆不同的容器写
同一套变换函数''逼出来的；第 9 章的单子是从''三种不同的效应居然共享同一套复合规则''
逼出来的。

\section{怎么读}\label{ux600eux4e48ux8bfb}

\begin{itemize}
\tightlist
\item
  \textbf{数学恐惧症}：标 ◇ 的段落都可以跳过，只读定义和代码。全书正文只保留''足以看懂代码''
  的数学。
\item
  \textbf{只想学 Haskell}：第 1--4 章、第 11--14 章是一条完整的工程路径，中间三部分当参考手册。
\item
  \textbf{只想学范畴论}：第 5、7、8、9 章是主线，代码当例子看。
\end{itemize}

\section{约定}\label{ux7ea6ux5b9a}

\begin{itemize}
\tightlist
\item
  代码全部来自 \passthrough{\lstinline!code/!} 里可编译的工程，正文里的每一段都能在仓库里找到并跑通。
\item
  术语首次出现给英文原文，全书中英对照见附录 C。
\item
  版本：GHC 9.14 / cabal 3.18，见附录 D。
\end{itemize}

本书在线版的源码在 \url{https://github.com/Freakwill/haskell-category}，
勘误与新章节都在那里。
